\documentclass[11pt]{article}
\usepackage{amsmath,amssymb,amsthm}
\newtheorem{theorem}{Theorem}\newtheorem{proposition}{Proposition}
\title{Pure loop sliding with labelled invocation paths: a directed native law}
\date{16 November 2026}
\begin{document}\maketitle
\paragraph{Audit and changed interface.} Separate next-delta manuscript. The no-initial-loop-port theorem, reverse obstructions and final well-founded-domain assembly section were independently reviewed within their stated scope. The source is the global-identifier incremental arena doctrine and its explicitly loop-labelled guarded compiler. This proves a directed pure-interface law, not JSV dinaturality or an invertible full-grade equation. All configurations, zero outcomes, full unused weighted polynomial grades and actual coefficient arrows remain present. The operational configuration inverse below is not generally a native inverse: global identifier and grade maps need not invert.

\section{Typed comparison}
Let $g:V\to U$ be a set function and let $P:A+U\rightsquigarrow B+V$ be an incremental arena with full weighted polynomial grade $C$. Assume no designated initial configuration of $P$ emits a $V$ port. Input reindexing retains original identifiers, rather than tagging them by the new input. Output transport retains original configurations and merely changes their displayed output labels. Put
\[
 P_V=(1_A+g)^*P:A+V\rightsquigarrow B+V,
 \qquad P_U=(1_B+g)_*P:A+U\rightsquigarrow B+U.
\]
Both are guarded. Their full grades are $C$, not input-dependent selected subcollections.
\begin{theorem}
There is a native arena map
\[
 S_{g,P}:\operatorname{Fb}_V(P_V)\longrightarrow
          \operatorname{Fb}_U(P_U).\tag{S}
\]
Its configuration function is bijective, its enabled-event and output-port maps are bijections at each configuration, and its primitive raw/coefficient maps are identities. Its full grade is the all-shape map
\[
 R_g C:R_V C\longrightarrow R_U C
\]
that applies $g$ to every internal edge label and leaves all node shapes and actual content factors unchanged. The map need not have any native inverse.
\end{theorem}
\begin{proof}
A source configuration is a finite invocation tree of original $P$ configurations. The root is over $a\in A$; a child attached to an original port $p$ carrying $v\in V$ has an input-reindexed configuration $(v,s)$, where the original configuration $s$ has input $g(v)$. Delete the redundant reindex tag at each child and replace its invocation edge identifier $(p,v)$ by $(p,g(v))$. The result is exactly a target compatible invocation tree. Conversely a target parent retains its original $P$ configuration, not only its transported output label. Its original label $v$ at each recursive port therefore determines the unique reindex tag for its child. This reconstructs the source tree even if $g$ is noninjective or nonsurjective. The argument applies to every compatible configuration, not just reachable ones.

On global identifier occurrence unions, apply $(p,v)\mapsto(p,g(v))$ letterwise to invocation words and keep the final original event or port identifier. This is one globally defined function. Within a compared configuration different simultaneous ports remain different, and the reconstruction gives the claimed local bijections. Initials, new-child initialization, old-child persistence, ordered output rendering and all transition diamonds commute. Enabled rows and their coefficient objects coincide, so the actual row arrows are identities. The terminal full body square is precisely $R_g C$: every completed invocation has the same original terminal shape and ordered ports. The content arrow is identity on each individually labelled node factor, with only the declared tensor coherence maps. This proves the native arrow without asserting a global inverse identifier function.
\end{proof}

\paragraph{Controls and independent evaluation.}
The configuration reconstruction identifies every legal structural history and every conditional coupling port. A policy includes the selected scope and its whole joint law at every such history, including zero and off-policy histories. Transporting these tables along the displayed history bijection is invertible: no conditional averaging, forward policy descent, or change of schedule is involved. This is an operational statement about tables, not an inverse full native arrow. The canonical contravariant control lift along (S) is exactly this transport because all local raw maps and marginals agree. The selected terminal evaluation has the same complete raw histories, weights and actual tensor coefficients, with the directed full-grade/body map $R_g C$. Independent terminal evaluation is the corresponding native map in the countable subprobability process category. No strict prefix clock identification beyond this same-step operational comparison is needed.

\section{Coherence}
For $g:V\to U$ and $h:U\to W$, use one original arena $P:A+W\rightsquigarrow B+V$. First slide $g$ through $(1_A+h)^*P$, then slide $h$ through $(1_B+g)_*P$. After the canonical commuting input-reindex/output-transport identification, the composite equals sliding $hg$ directly. Indeed both delete the same redundant input tags, send each invocation edge $(p,v)$ to $(p,hg(v))$, and retain every original event/raw/coefficient datum. On every unused grade shape the maps are
\[
 R_h C\,R_g C=R_{hg}C.
\]
Identity sliding is the identity under the structural input-reindex identification. For a native map $P\to P'$ at the fixed endpoints $A+U,B+V$, sliding is natural with the same fixed $g$: both routes apply the supplied original identifier functions and the same edge-label function, and their actual node coefficient arrows agree. No additional changing-endpoint naturality is asserted without specifying its commuting endpoint squares. These are directed pure-loop functoriality laws, not an unrestricted contraction or feedback interchange assertion.

\section{Two independent reverse obstructions}
\begin{proposition}[Unused-shape obstruction]
A reverse native comparison need not exist even when every realized recursive tree is trivial.
\end{proposition}
Take $V=\varnothing$, $U=\{u\}$, and $C$ with one binary shape and unit coefficient. An original event-free $P:A+U\rightsquigarrow B$ can emit two $B$ ports. The source full grade $R_\varnothing C$ has only arity-two shapes. The target $R_U C$ also contains the formal two-node tree with one internal edge and three external leaves. No arity-preserving map from the full target grade to the source grade exists. The unused tree cannot be discarded merely because this particular body never emits a loop output.

\begin{proposition}[Positive-body obstruction for a surjective loop map]
A reverse native comparison can fail for surjective $g$, even though all shapes in both grades have the same arity and the loop map admits a set-theoretic section.
\end{proposition}
Let $A=\{a\}$, $B=\{b\}$, $V=\{v_0,v_1\}$, $U=\{u\}$, and $g(v_i)=u$. Let $C$ have two unary shapes $c_R,c_T$, both with unit content coefficient. At input $a$, the atomic arena $P$ draws a fair bit $i$, emits its single port with label $v_i$, and has terminal shape $c_R$. At input $u$, it deterministically emits $b$ with shape $c_T$. All event coefficient objects are the tensor unit. Initial configurations emit no ports, so the guard holds.

The source feedback has exactly two complete raw traces of weight $1/2$, with distinct terminal full shapes
\[
 t_0=(c_R\xrightarrow{v_0}c_T),\qquad
 t_1=(c_R\xrightarrow{v_1}c_T).
\]
The target also has two complete raw traces of weight $1/2$, but both have the same full shape $t=(c_R\xrightarrow{u}c_T)$. All other formal trees and their coefficients remain in both grades. Suppose any reverse native process map existed between the independent terminal evaluations, allowing an arbitrary arity-preserving full grade map $\rho$. The body square forces the image of every positive target raw trace to have source shape $\rho(t)$. Exact probability pushforward requires positive mass at each of the two source traces, whose shapes differ. This is impossible. Therefore no such native evaluated-process map exists, and functorial independent evaluation excludes a reverse native arena map as well. No assumption that the reverse raw map or grade map is the evident inverse has been made.

The distinction is material: the forward comparison preserves the complete raw bit, whereas its grade map identifies two positive shape labels. Configurations can still recover the bit from the original parent state. Recoverability at that level does not supply a single reverse grade function compatible with both bodies. A section of $g$ chooses one loop label, not both simultaneous positive body obligations.

\section{Assembly with well-founded silent initialization}
Drop the earlier no-initial-$V$-port hypothesis on $P$. Keep finite initial
port sets and the same labelled compiler, now with finite silent initial
closures. Define the following two independently specified relations:
\[
 v\mathrel{S}w\iff i_P(g(v))\text{ emits a port originally labelled }w\in V,
 \qquad
 u\mathrel{T}u'\iff \exists w\in V\;
 [i_P(u)\text{ emits }w\text{ and }g(w)=u'].
\]
These are exactly the initial-loop relations of $P_V$ and $P_U$, respectively.
Repeated port occurrences remain separate in actual trees, though the
relations only record existence.
\begin{theorem}[Exact common admission and silent sliding]
The relation $S$ is well founded if and only if $T$ is well founded.
On this common domain, (S) extends to an actual native comparison of the
well-founded-silent feedback arenas, with the same full map $R_g C$,
the same operational configuration/history bijections and the same actual
primitive coefficient arrows. It transports all complete policy tables and
preserves paid event count. Identity, composition and fixed-endpoint native
naturality proved above remain valid whenever the displayed feedbacks are
admitted. Neither surjectivity nor injectivity of $g$ is required.
\end{theorem}
\begin{proof}
An infinite $S$ chain $v_0 S v_1 S\cdots$ pushes forward to the infinite
$T$ chain $g(v_0)Tg(v_1)T\cdots$. Conversely suppose $u_0Tu_1T\cdots$.
For each step choose one original emitted label $w_i$ at $i_P(u_i)$ with
$g(w_i)=u_{i+1}$. Then $w_i S w_{i+1}$: the latter label is emitted by
$i_P(u_{i+1})=i_P(g(w_i))$. Thus the target chain, after its first step,
lifts to an infinite source chain. This uses the original emitted labels,
not a section of $g$, and covers target starting inputs outside its image.
The usual dependent-choice setting is the same as in the finite-branching
silent-domain theorem. Finite branching of the original initial port sets
now makes both initial closures finite at every loop input and every
external input $a$.

The sets of all compatible finite configurations are those of the earlier
sliding construction, irrespective of which are designated initial.
Its map deletes child reindex tags; its inverse reads each original
parent port label. Apply this reconstruction to the finite initial trees.
The root original initial state is the same; every child has the original
initial state at $g(v)$ on both sides. Finite tree recursion therefore
identifies the two designated initial closures. At a paid event, each newly
born original $V$ port attaches precisely these corresponding complete
initial closures. No paid event is executed by the attachment. Old children,
all original event rows, their actual coefficient objects, and all
persistence diamonds are unchanged. The labelled global identifier maps
are exactly the previous letterwise functions $(p,v)\mapsto(p,g(v))$.
Consequently the same native proof and terminal $R_g C$ body square apply.
Silent invocation nodes still contribute their full content factors; they
introduce no event context factor and are not removed from any grade.

History reconstruction, including every zero and off-policy history,
now starts with the corresponding finite initial trees and commutes with
each paid transition. The transported scope contains the same primitive
events and the same conditional law. Thus every policy table transports
as before, with unchanged paid event counts. Identity and composite label
maps act by the same word functions on these larger initial trees. The
previous coherence proof therefore applies without quotienting any silent
node or unused grade shape.
\end{proof}

For fixed $a,n$, let $\mathcal K_V(a,n)$ and $\mathcal K_U(a,n)$ be the
finite structural sets reached by at most $n$ singleton paid events,
including every raw outcome, from the respective silent initial closures.
The preceding transition bijection identifies these sets and preserves
the whole invocation-node tree. Hence in this particular comparison the
structural depth bounds agree exactly:
\[
 N_V(a,n)=N_U(a,n).
\]
Depth counts invocation nodes, so the equality includes zero-cost trees of
depth greater than one. This equality is special to sliding's operational
bijection; no such equality is inferred for arbitrary native maps.
Under transported policies, the same-step absorbing prefix histories are
identified. Their terminal domains correspond exactly; active histories
receive no invented terminal shape. At every first-hit cost layer, including
cost zero, the raw witnesses and event tensors agree and the full body map
is $R_g C$. These are actual forward native maps of the fixed-grade
partial-terminal-body towers and terminating cost layers. The independent
macro cost construction from the silent-domain theorem uses the equal
bounds above and the corresponding complete local traces, so its comparison
with the event-cost construction commutes with sliding on raw trees and
actual tensor/content arrows. When $A,B,U,V$ are countable, the cited countable matching-trace
interface yields the same directed terminal comparison. The finite initial,
prefix and fixed-cost statements above require no countability of these
endpoint sets. This is not an inverse grade map or an
identification with an unrelated instantaneous matching trace.

\end{document}
